\documentclass[10pt,a4paper]{amsart}
\usepackage[margin=19mm]{geometry}
\usepackage{amsmath,amssymb,mathtools}
\usepackage[scr=rsfs,cal=euler]{mathalfa}
\usepackage{xfrac}
\usepackage[unicode,hidelinks,bookmarks=false]{hyperref}
\newcommand{\ie}{i.e.\ }
\newcommand{\cf}{cf.\ }
\newcommand{\resp}{respectively\ }
\newcommand{\Subsection}{\subsection}
\providecommand{\nobreakdash}{\nobreak\hbox{-}}
\setlength{\emergencystretch}{2em}
\multlinegap0pt
\usepackage{amssymb}
\usepackage{mathtools}
\usepackage{bm}
\newtheorem{theorem}{Theorem}
\title{Multidimensional scalar conservation laws with discontinuous flux: well-posedness without non-degeneracy}
\author{Darko Mitrovi\'c}
\date{Source: arXiv:2607.17310v1 (CC BY 4.0)}
\begin{document}
\maketitle

\noindent\emph{Source and licence.} Multidimensional scalar conservation laws with discontinuous flux: well-posedness without non-degeneracy. Version arXiv:2607.17310v1 (CC BY 4.0), distributed by arXiv under the Creative Commons Attribution 4.0 International licence, http://creativecommons.org/licenses/by/4.0/, as recorded in the arXiv licence metadata for this exact version and shown on the abstract page https://arxiv.org/abs/2607.17310v1. Full licence notice: \texttt{licenses/arxiv-2607.17310-CC-BY-4.0.txt}.\par\smallskip

\noindent\textit{Editorial scope.} Counted unit: the article's theorem thm:strong-trace-normal-flux-quotient (source0.tex lines 1795--1884). The article's complete original proof of it is delivered with it (source0.tex lines 1886--2069, one \texttt{proof} environment of 184 lines). No mathematics is written, repaired or re-derived: the statement and the proof are the article's own lines, in the article's own notation, and no retained line is substituted or re-flowed.  No further source-local context is needed: every cross-reference of every retained line resolves inside the two delivered spans.  Nothing was omitted from any retained span, so the package carries no omission record.  No retained line contains a citation, so the wrapper carries no bibliography and no bibliography entry.  No retained line contains a figure environment, an image-inclusion command or drawing code, so no image is delivered and the package declares no asset dependency.  Editorial formatting only, in addition: an AMS \texttt{amsart} preamble that restates the article's own declarations and macros used by the retained lines, namely the environments \texttt{theorem} on separate counters, together with the standard packages those lines need.  The retained environments print in this document as Theorem 1 in order of appearance, whereas the article prints the same items with the numbering of its own full document.  The complete native source of the article as distributed in the arXiv e-print for this exact version is bundled unchanged as \texttt{sources/205554/source0.tex}.
\begin{theorem}[Strong trace of a normal-flux quotient limit]
\label{thm:strong-trace-normal-flux-quotient}
Let
\[
    Q_R^+
    :=
    B_R'\times(0,R),
    \qquad
    x=(x',y),
    \qquad
    K=[-a,a],
\]
and assume that
\[
    N\in C^1\bigl(\overline{Q_R^+}\times K\bigr).
\]
Set
\[
    D_N(x,\lambda)
    :=
    \partial_\lambda N(x,\lambda)
\]
and define the normal-flux quotient by
\begin{equation}
\label{eq:normal-flux-quotient-trace-theorem}
    \Pi_N(x,\lambda)
    :=
    \int_{-a}^{\lambda}
        |D_N(x,s)|^2\,ds.
\end{equation}

The boundary non-degenerate region is
\begin{equation}
\label{eq:boundary-normal-nondegenerate-region}
    \widetilde A_N
    :=
    \left\{
        (x',\lambda)\in B_R'\times K:
        D_N(x',0,\lambda)\neq0
    \right\}.
\end{equation}

Let
\[
    u_n:Q_R^+\to K
\]
be measurable, and assume that
\begin{equation}
\label{eq:normal-quotient-interior-limit}
    \Pi_N(x,u_n(x))
    \longrightarrow z(x)
    \qquad
    \text{strongly in }L^1_{\mathrm{loc}}(Q_R^+).
\end{equation}

Assume moreover that, for every cylinder
\[
    B\times I\Subset\widetilde A_N,
\]
there exists \(\delta>0\) such that, for every
\[
    [\alpha,\beta]\Subset I,
\]
the constant truncations
\[
    s_{\alpha,\beta}(u_n)
\]
satisfy Panov's compactness criterion on
\[
    B\times(0,\delta),
\]
and their entropy productions are locally uniformly bounded up to
\[
    B\times\{0\}.
\]

Then \(z\) has a strong trace
\[
    z^\tau\in L^\infty_{\mathrm{loc}}(B_R')
\]
on \(B_R'\times\{0\}\).  More precisely, for every
\(B\Subset B_R'\),
\begin{equation}
\label{eq:normal-quotient-strong-trace}
    \operatorname*{ess\,lim}_{y\to0+}
    \int_B
        |z(x',y)-z^\tau(x')|\,dx'
    =0.
\end{equation}
\end{theorem}

\begin{proof}
The set \(\widetilde A_N\) is relatively open.  Let
\[
    B\times I\Subset\widetilde A_N.
\]
Then
\[
    |D_N(x',0,\lambda)|\ge c>0
    \qquad
    \text{on }\overline B\times\overline I.
\]
By continuity, after decreasing \(\delta>0\) if necessary,
\[
    D_N(x',y,\lambda)\neq0
\]
on
\[
    \overline B\times[0,\delta]\times\overline I.
\]
Consequently, the canonical polynomial vector
\[
    \mathcal N_N(x,\lambda)
    =
    \bigl(
        N(x,\lambda)^2,\,
        N(x,\lambda),\,
        N(x,\lambda)^3,\ldots,
        N(x,\lambda)^m
    \bigr)
\]
satisfies Panov's interval non-degeneracy condition on this cylinder.

By Panov's compactness theorem and a diagonal extraction, the constant
truncations
\[
    s_{\alpha,\beta}(u_n),
    \qquad
    [\alpha,\beta]\Subset I,
\]
converge locally to quasi-solutions on
\[
    B\times(0,\delta).
\]
The assumed entropy-production bounds permit the application of Panov's trace
theorem.  Hence the limiting constant truncations have strong traces on
\(B\times\{0\}\).

For a rectangular weight
\[
    \rho(x',\lambda)
    =
    \psi(x')\mathbf 1_{(\alpha,\beta)}(\lambda),
    \qquad
    \operatorname{supp}\psi
    \times[\alpha,\beta]
    \Subset\widetilde A_N,
\]
one has
\[
\begin{aligned}
    \int_K
        \operatorname{sgn}(u_n-\lambda)
        \rho(x',\lambda)\,d\lambda
    &=
    \psi(x')
    \bigl(
        2s_{\alpha,\beta}(u_n)-\alpha-\beta
    \bigr).
\end{aligned}
\]
Therefore the limit of this average has a strong trace.  Approximation by
finite sums of rectangular weights gives the same conclusion for every
continuous weight compactly supported in \(\widetilde A_N\).

We apply this conclusion to the boundary weight
\[
    \rho_N^\tau(x',\lambda)
    :=
    |D_N(x',0,\lambda)|^2.
\]
This weight vanishes outside \(\widetilde A_N\).  On each compact boundary
patch it can be approximated uniformly by continuous weights compactly
supported in \(\widetilde A_N\).  For example, one may use
\[
    \rho_{N,j}^\tau(x',\lambda)
    :=
    |D_N(x',0,\lambda)|^2
    \vartheta_j\bigl(|D_N(x',0,\lambda)|\bigr),
\]
where
\[
    \vartheta_j(r)=0
    \quad\text{for }r\le\frac1j,
    \qquad
    \vartheta_j(r)=1
    \quad\text{for }r\ge\frac2j.
\]
The discarded contribution satisfies
\[
\begin{aligned}
&\left|
    \int_K
        \operatorname{sgn}(u_n-\lambda)
        \bigl(
            \rho_N^\tau-\rho_{N,j}^\tau
        \bigr)(x',\lambda)\,d\lambda
\right|
\\
&\qquad\le
    \frac{4|K|}{j^2},
\end{aligned}
\]
uniformly in \(n\).  Hence the limit of the averages with weight
\(\rho_N^\tau\) has a strong trace.

Since \(N\in C^1(\overline{Q_R^+}\times K)\),
\[
    |D_N(x',y,\lambda)|^2
    \longrightarrow
    |D_N(x',0,\lambda)|^2
\]
locally uniformly in \((x',\lambda)\) as \(y\to0+\).  Therefore the same
trace is obtained for the limit of
\[
    G_n(x)
    :=
    \int_K
        \operatorname{sgn}(u_n(x)-\lambda)
        |D_N(x,\lambda)|^2\,d\lambda.
\]

Finally, define
\[
    C_N(x)
    :=
    \int_K
        |D_N(x,\lambda)|^2\,d\lambda.
\]
For every \(n\),
\begin{equation}
\label{eq:normal-quotient-average-trace-identity}
    \Pi_N(x,u_n(x))
    =
    \frac12G_n(x)
    +
    \frac12C_N(x).
\end{equation}
By \eqref{eq:normal-quotient-interior-limit},
\[
    G_n
    \longrightarrow
    G:=2z-C_N
    \qquad
    \text{strongly in }L^1_{\mathrm{loc}}(Q_R^+).
\]
The preceding argument shows that \(G\) has a strong trace \(G^\tau\).

The function \(C_N\) is continuous up to the boundary, with
\[
    C_N^\tau(x')
    =
    \int_K
        |D_N(x',0,\lambda)|^2\,d\lambda.
\]
Hence
\[
    z
    =
    \frac12(G+C_N)
\]
has the strong trace
\begin{equation}
\label{eq:normal-quotient-trace-formula}
    z^\tau(x')
    :=
    \frac12
    \left[
        G^\tau(x')
        +
        \int_K
            |D_N(x',0,\lambda)|^2\,d\lambda
    \right].
\end{equation}
\end{proof}

\end{document}
